\documentclass[11pt]{article}
\usepackage[margin=1in]{geometry}
\usepackage[T1]{fontenc}
\usepackage{lmodern}
\usepackage[hidelinks]{hyperref}
\hypersetup{pdftitle={SimTIO: Cover letter},pdfauthor={Ruimin Ke}}
\setlength{\parindent}{0pt}
\setlength{\parskip}{0.8em}
\pagestyle{empty}
\begin{document}
Dear Editors of \textit{Multimodal Transportation},

Please consider our manuscript, ``SimTIO: A Simulation-Grounded Multi-Agent LLM Framework for Compositional Traffic Intervention Optimization,'' for publication as a Full Length Article.

The study addresses how language-model proposals can be translated into executable traffic interventions and assessed against measured network effects. SimTIO combines role-specific specialists, program-generated legal mutation catalogs, executable validation, SUMO evaluation, and explicit acceptance constraints. Compatible interventions can be composed and refined using feedback from their joint simulation, while a no-operation option remains available.

The contribution is an auditable scenario-screening procedure and a controlled assessment of the language model's role. A frozen-input comparison holds initial candidates, parents, mutation catalogs, simulator seeds, and search rules constant while varying the selector. A broader adaptive search evaluates five initial candidates, two children per selected parent, and continued joint-feedback refinement. Additional experiments assess sampling variability, synthetic demand intensity, acceptance thresholds, and spatially concentrated adverse effects.

The results do not establish an LLM advantage over random or deterministic selection. They show why matched search structures and explicit spatial diagnostics matter when evaluating LLM-based transportation decision support. All experiments use uncalibrated synthetic passenger demand; the manuscript does not claim demonstrated real-city improvements or validated multimodal benefits.

We believe the manuscript is relevant to the journal's interests in intelligent transportation, emerging computational methods, and transport network modeling. It distinguishes operational interventions from planning scenarios and discusses the evidence required before practical deployment.

Thank you for considering our manuscript.

Sincerely,\\
Ruimin Ke\\
Corresponding author, on behalf of Shuyang Li and Ruimin Ke\\
Civil and Environmental Engineering\\
Rensselaer Polytechnic Institute\\
110 8th Street, Troy, New York 12180-3590, USA\\
\href{mailto:ker@rpi.edu}{ker@rpi.edu}
\end{document}
